\documentclass[11pt]{article}
\usepackage[T1]{fontenc}
\usepackage{lmodern,amsmath,amssymb,amsthm,booktabs,array,microtype,xcolor}
\usepackage[letterpaper,margin=1in]{geometry}
\usepackage[colorlinks=true,linkcolor=blue!45!black,urlcolor=blue!45!black]{hyperref}
\newtheorem{theorem}{Theorem}[section]
\newtheorem{lemma}[theorem]{Lemma}
\newtheorem{proposition}[theorem]{Proposition}
\newtheorem{corollary}[theorem]{Corollary}
\newtheorem{remark}[theorem]{Remark}
\setlength{\emergencystretch}{3em}
\title{Additive stability for terminal path counts}
\author{Collaborative research working draft}
\date{10 November 2026}
\begin{document}
\maketitle
\begin{abstract}
For every fixed $\varepsilon>0$, we prove that a connected graph of cycle rank $b$ with at least $(1/2+\varepsilon)2^b$ simple paths between two specified terminals has terminal distance at least $b-C_\varepsilon$.
The constant depends only on $\varepsilon$.
The reduced terminal kernel also has bounded pathwidth, and deleting a suitable shortest path from its cubic expansion leaves bounded cycle rank.
The proof combines the source manuscript's cycle-space machinery, an extended fixed-gap forest estimate, Fable 5.2's localization and distance-layer argument, and a new short-cycle packing argument in bounded layers.
The proof has received independent internal audits; it has not been formalized in Lean.
\end{abstract}

\section*{Scope and attribution}
This note uses the forest-probability machinery of the November 10, 2026 manuscript
\emph{The number of cycles in a graph and the minimum number of edges in a pancyclic graph},
whose local source is \nolinkurl{paper/cycle-count/pancyclic-paper.tex} at revision \texttt{a965788}.
The source manuscript's cut, spectral, and second-moment lemmas are inputs to the parameter extension below.
They are not re-proved here.
Fable 5.2 supplied localization and the shortest-path layer proof of general linear stability.
The Codex path-stability agent supplied the deduction of additive stability from bounded layers.
The Codex root and four-regular agents independently audited that deduction.
The original manuscripts and formal sources have not been changed.

Rank means $\beta(G)=|E(G)|-|V(G)|+\operatorname{comp}(G)$.
The forest probability $p(G,W)$ is the probability that the part outside $E(W)$ of a uniform even edge set is a forest.
The main conclusion is Theorem~\ref{thm:general-additive-path-stability}.
Earlier estimates are retained because their proofs supply its inputs.

\tableofcontents
\section{Path stability and the size of a marked subgraph}
\label{sec:research-path}

These results use the forest proof in the November 10 manuscript.
The parameter extension below is a mathematical argument, not a Lean formalization.
It first improves the general path bound to order $b/\log b$.
Theorem~\ref{thm:general-linear-path-stability} later obtains the full linear bound from this enlarged forest range.
Theorem~\ref{thm:general-additive-path-stability} then strengthens it to an additive loss depending only on the density gap.
For two-terminal series-parallel networks, an elementary argument gives an explicit logarithmic loss.

All logarithms in this section have base two, except for $\ln$ and $\exp$.
Write $P_{uv}(N)$ for the number of simple paths between distinct vertices $u,v$.
Write $b=\beta(N)$.

\subsection{From paths to forest outside parts}

\begin{lemma}\label{lem:research-path-injection}
Let $N$ be a connected simple graph, and let $A$ be any $u$--$v$ path of length $\ell$.
Add a new vertex $w$ and edges $uw,wv$ to obtain $H$.
Let $W=A\cup\{uw,wv\}$.
Then
\[
 \beta(H)=b+1,\qquad \beta(W)=1,\qquad |E(W)|=\ell+2,
 \qquad P_{uv}(N)\le p(H,W)2^b.
\]
\end{lemma}

\begin{proof}
Each terminal path $P$ gives the cycle $P\cup\{uw,wv\}$ of $H$.
Its outside part is $P\setminus E(A)$, which is a forest.
Suppose that $P\setminus E(A)=Q\setminus E(A)$.
Then $P\mathbin\triangle Q$ is an even edge set supported on the forest $A$.
It must be empty, so $P=Q$.
There are $2^{\beta(H)-\beta(W)}=2^b$ outside parts.
Their forest proportion is $p(H,W)$.
\end{proof}

The formalized forest theorem gives $\ell>\sqrt{b+1}-2$ for fixed positive path-density gap and sufficiently large $b$.
The manuscript's size extension gives $\ell>(b+1)/\log^2(b+1)-2$.
The next argument removes one logarithmic factor for a fixed gap.

\subsection{A fixed-gap forest estimate at size \texorpdfstring{$r/\log r$}{r/log r}}

\begin{theorem}\label{thm:research-marked-size}
Fix $0<\delta<1/2$, and set
\[
 A_\delta=6\exp(16/\delta),\qquad
 a_\delta=\frac{1}{12A_\delta}.
\]
For sufficiently large $r$, the following statement holds.
Let $G$ be a connected simple graph of rank $r$.
Let $W$ be a nonempty union of cycles satisfying
\[
 |E(W)|\le a_\delta\frac{r}{\log r},\qquad
 \operatorname{comp}(W)\le\frac{\log r}{100}.
\]
Then $p(G,W)\le1/2+\delta$.
\end{theorem}

\begin{proof}
We use the manuscript's reduction, cut lemmas, spectral estimate, and second-moment argument.
The details below check each estimate affected by the new parameters.
Constants in $O_\delta(\cdot)$ may depend on $\delta$.

Put $x=\log r$, $y=\log x$, $D=\lfloor x/10\rfloor$, and $s=\lfloor(D-1)/2\rfloor$.
Set $L=\lfloor A_\delta x\rfloor$.
The reduction gives a subcubic multigraph $\Gamma$ and a marked vertex set $P$.
Every vertex outside $P$ has degree three, $|P|\le2|E(W)|$, and $\Gamma[P]$ has at most $x/100$ components.
Also $\beta(\Gamma)=r$ and $2r-2\le |V(\Gamma)|\le2r+2|E(W)|$.
Let $\mathcal F$ mean that the random even subgraph is a forest on $\Gamma-P$.
It suffices to rule out $\mathbb P(\mathcal F)>1/2+\delta$.

Assume this strict inequality.
Choose a maximal vertex-disjoint family $\mathcal Q$ of cycles of length at most $D$ in $\Gamma-P$.
The small-cut lemma gives
\[
 |\mathcal Q|<\frac{(D+1)2^{D-1}}{\delta}
       =O_\delta(xr^{1/10}).
\]
For their vertex union $S$, this implies
\[
 |S|=O_\delta(x^2r^{1/10})=o(r/x).
\]
Let $J=\Gamma-(P\cup S)$, and put $v=|V(J)|$.
This graph is simple and has girth greater than $D$.
The vertex and deficit estimates become
\[
 v=2r+O_\delta(r/x),\qquad
 d_J:=\sum_{z\in V(J)}(3-\deg_Jz)
 \le 3(|P|+|S|)
 \le\frac{6a_\delta r}{x}+o(r/x).
\]
Consequently,
\[
 \frac{Ld_J}{2v}\le\frac32 A_\delta a_\delta+o(1)
                 =\frac18+o(1).
\]

Let $B$ be the non-backtracking matrix of $J$.
For even $j$ with $6x\le j\le L$, the spectral estimate gives
\[
 2^{-j}\operatorname{tr}B^j
 \ge\left(1-\frac{d_J}{2v}\right)^j-3v2^{-j/2}
 \ge1-\frac{Ld_J}{2v}-3vr^{-3}\ge\frac12
\]
for sufficiently large $r$.
For the cycles $C$ of $J$ of length at most $L$, put
$\mu=\sum_C2^{-|C|}$.
The repeated-vertex walk estimate from the manuscript has error at most
\[
 \frac32 L^2 2^{-s}\mu=o(\mu).
\]
Therefore
\[
 (1+o(1))\mu
 \ge\frac12\sum_{\substack{6x\le j\le L\\j\ \mathrm{even}}}\frac1{2j}
 =\frac18\ln(A_\delta/6)+o(1).
\]
In particular, $\mu\ge\frac1{16}\ln(A_\delta/6)=1/\delta$ for sufficiently large $r$.

It remains to check the pair correlations.
Use $t_0=\lfloor y/2\rfloor$ in the manuscript's link argument.
For disjoint cycles $C,C'$ of length at most $L$, it gives
\[
 \ell(C,C')\le \operatorname{comp}(\Gamma[P])
 +\frac{2L}{t_0}
 +\frac{(t_0+1)2^{t_0-1}}{\delta}
 +\left(\frac{8L}{D}\right)^2
 \le\frac{x}{100}+o(x).
\]
Here the three added terms are respectively $O_\delta(x/y)$, $O_\delta(y\sqrt{x})$, and $O_\delta(1)$.
The exceptional-partner estimate thus remains
\[
 2^{x/100+o(x)}L^3 2^{-s}
 \le2^{x/100-x/20+o(x)}=o(1).
\]
The same weighted variable $Z$ therefore satisfies
\[
 \mathbb EZ=\mu,\qquad
 \mathbb EZ^2\le2\mu^2+(1+o(1))\mu.
\]
Since $\mathcal F$ implies $Z=0$, the second-moment lemma yields
\[
 \mathbb P(\mathcal F)\le\frac12+\frac{1+o(1)}{4\mu}
 \le\frac12+\frac\delta2
\]
for sufficiently large $r$.
This contradicts the assumption.
\end{proof}

\begin{corollary}\label{cor:research-general-path}
Fix $0<\varepsilon\le1/2$.
For sufficiently large $b$, if
\[
 P_{uv}(N)\ge(1/2+\varepsilon)2^b,
\]
then every terminal path has length
\[
 \ell>a_{\varepsilon/2}\frac{b+1}{\log(b+1)}-2.
\]
In particular, $\ell=\Omega_\varepsilon(b/\log b)$.
\end{corollary}

\begin{proof}
Apply Lemma~\ref{lem:research-path-injection}.
If the displayed length inequality failed, Theorem~\ref{thm:research-marked-size} would give
$P_{uv}(N)\le(1/2+\varepsilon/2)2^b$, a contradiction.
The one-component condition holds for sufficiently large $b$.
\end{proof}

The threshold is not explicit because the supporting asymptotic estimates are not explicit.
The additive $-2$ comes from the added two-edge path.
Neither issue affects the asymptotic order.

\begin{corollary}\label{cor:research-small-marked}
Suppose $r\to\infty$ along connected simple graphs $G$ and nonempty unions of cycles $W$.
If $|E(W)|=o(r/\log r)$ and $\operatorname{comp}(W)\le(\log r)/100$, then
$p(G,W)\le1/2+o(1)$.
\end{corollary}

\begin{proof}
For each fixed $0<\delta<1/2$, the size hypothesis eventually meets Theorem~\ref{thm:research-marked-size}.
Thus $\limsup p(G,W)\le1/2$.
\end{proof}

\subsection{An additive theorem for two-terminal series-parallel networks}

A two-terminal series-parallel network starts with one edge and its two endpoints.
Series composition identifies one terminal from each network.
Parallel composition identifies the corresponding terminal pairs.
Apart from those identifications, the two networks are disjoint.
Parallel edges are permitted and are counted separately.
This definition imposes a condition on the terminals, not only on the underlying graph.

\begin{theorem}\label{thm:research-ttsp}
Let $(N,u,v)$ be a two-terminal series-parallel network of rank $b$ and terminal distance $\ell$.
Then
\[
 P_{uv}(N)\le 2^{b-1}+2^{\ell-1}.
\]
Consequently, for $0<\varepsilon\le1/2$,
\[
 P_{uv}(N)\ge(1/2+\varepsilon)2^b
 \quad\Longrightarrow\quad
 \ell\ge b-\log\frac1{2\varepsilon}
 \quad\Longrightarrow\quad
 \ell\ge\frac{b}{\log(1/\varepsilon)}.
\]
\end{theorem}

\begin{proof}
Every terminal path is a join, so $P_{uv}(N)\le2^b$ for every connected network.
Define $h(t)=(1+2^{-t})/2$ for $t\ge0$.
We prove the stronger normalized statement
\[
 \frac{P_{uv}(N)}{2^b}\le h((b-\ell)_+),
 \qquad (z)_+=\max\{z,0\}.
\]
It holds for a single edge.

For series composition, write $b_i,\ell_i,P_i$ for the two input networks.
Then $b=b_1+b_2$, $\ell=\ell_1+\ell_2$, and $P=P_1P_2$.
For $s,t\ge0$,
\[
 h(s+t)-h(s)h(t)=\frac{(1-2^{-s})(1-2^{-t})}{4}\ge0.
\]
Also $(b_1-\ell_1)_++(b_2-\ell_2)_+\ge(b-\ell)_+$, and $h$ is decreasing.
The normalized statement follows by multiplication.

For parallel composition, $b=b_1+b_2+1$, $\ell=\min\{\ell_1,\ell_2\}\ge1$, and $P=P_1+P_2$.
Since $b_1,b_2\ge0$,
\[
 P\le2^{b_1}+2^{b_2}
 \le2^{b_1+b_2}+1
 =2^{b-1}+1
 \le2^{b-1}+2^{\ell-1}.
\]
Together with $P\le2^b$, this is the normalized statement.
Induction proves the theorem's first inequality.

The density hypothesis now gives $\varepsilon2^b\le2^{\ell-1}$.
Thus $b\le\ell+\log(1/(2\varepsilon))$.
Since $\ell\ge1$, the right side is at most $\ell\log(1/\varepsilon)$.
\end{proof}

The coefficient is sharp when $\varepsilon=2^{-b}$.
Take a series of $b-1$ triangles between $u,v$, and add the edge $uv$.
For $b\ge3$, the result is simple, has rank $b$, distance one, and $2^{b-1}+1$ terminal paths.
For $b=2$, a direct edge in parallel with two internally disjoint two-edge paths has the same count.
For $b=1$, use a triangle with adjacent terminals.

The theorem fails for arbitrary choices of terminals in a series-parallel graph.
On vertices $0,1,2,3,4$, take
\[
 E=\{01,02,03,04,12,14,23\},\qquad u=3,\quad v=4.
\]
This graph has rank three and terminal distance two.
Its seven terminal paths are
\[
 3014,\quad30214,\quad304,\quad32014,\quad3204,\quad32104,\quad3214.
\]
Thus $P=7>2^{3-1}+2^{2-1}=6$.
The graph is a triangle with two vertices successively attached to existing edges, so it has no $K_4$ minor.
It does not satisfy the two-terminal composition hypothesis for terminals $3,4$.

\subsection{An elementary structural reduction}

\begin{proposition}\label{prop:research-offroute}
Let $N$ be connected, with distinct terminals $u,v$.
Let $N_0$ be the union of the blocks on the $u$--$v$ route in the block-cut tree.
Count bridges as blocks.
Put $s=\beta(N)-\beta(N_0)$.
Then
\[
 P_{uv}(N)=P_{uv}(N_0)\le2^{\beta(N_0)}=2^{b-s}.
\]
In particular, $P_{uv}(N)>2^{b-1}$ implies $s=0$.
\end{proposition}

\begin{proof}
A simple terminal path cannot enter a block outside the terminal route.
It would have to revisit the articulation vertex where it left that route.
Thus all terminal paths lie in $N_0$.
Apply the join bound to $N_0$.
Ranks add across blocks, so $s$ is a nonnegative integer.
\end{proof}

This removes cyclic blocks outside the terminal route from the unresolved problem.
It does not control the internal structure of a single block.
The fixed-gap argument still loses a factor $\log b$ there.
Its spectral estimate counts walks of length at least order $\log b$.
A marked set of order $b$ creates a constant degree deficit, whose effect accumulates along those walks.
The present retuning does not prove general linear stability.

\paragraph{Finite checks.}
The accompanying script checks 1,728 distinct composition states for networks with at most 12 edges.
It verifies the inequality with exact integers across 27,264 series or parallel transitions.
It also enumerates the seven counterexample paths and the sharp simple examples for ranks three through nine.
These checks do not verify the asymptotic estimates or replace their proofs.

\section{General linear path stability}
\label{sec:general-linear-path-stability}

Fable 5.2 proposed the localization and shortest-path layer argument in its November 10 research notes.
The Codex root and path-stability agents independently audited that argument.
The proof below includes the minor padding correction found in both audits.
It uses Theorem~\ref{thm:research-marked-size} to supply the required forest estimate.
Thus the conclusion is a mathematical theorem, rather than a conditional reduction.
It has not been formalized in Lean.

All logarithms in this section have base two, except $\ln$.
Put $\Phi(x)=x/\log^2x$ for $x\ge8$.
This function is increasing on its domain.

\begin{theorem}[General linear stability]\label{thm:general-linear-path-stability}
For every $0<\varepsilon\le1/2$, there is a constant $c_\varepsilon>0$ with the following property.
Let $N$ be a connected simple graph of rank $b$, with distinct terminals $u,v$.
If
\[
 P_{uv}(N)\ge(1/2+\varepsilon)2^b,
\]
then every $u$--$v$ path has length at least $c_\varepsilon b$.
In particular, $\operatorname{dist}_N(u,v)\ge c_\varepsilon b$.

An explicit expression in terms of the forest threshold is available.
Set $\delta=\varepsilon/2$, and let $R_\delta$ be an integer as in Lemma~\ref{lem:general-fe2} below.
Put
\[
 M=2R_\delta,\qquad
 \Theta=\max\left\{2^{36},\frac{(M+1)M2^{M+1}}{\delta}\right\}.
\]
One may take $c_\varepsilon=(\Theta+1)^{-2}$.
Every terminal path also satisfies the stronger inequality
\begin{equation}\label{eq:general-linear-strong}
 \ell>\frac{2b}{\Theta+1}-(\Theta+1).
\end{equation}
\end{theorem}

The formula determines a positive constant once the forest threshold is fixed.
The present proof does not give a numerical value for that threshold.

\subsection{The forest input and its localization}

\begin{lemma}\label{lem:general-fe2}
Fix $0<\delta<1/2$.
There is an integer $R_\delta\ge2^{100}$ such that the following statement holds.
If $G$ is connected and simple, $\rho=\beta(G)\ge R_\delta$, and $W$ is one cycle with
$|E(W)|\le\Phi(\rho)$, then
\[
 p(G,W)\le1/2+\delta.
\]
\end{lemma}

\begin{proof}
Use $a_\delta$ from Theorem~\ref{thm:research-marked-size}.
Choose $R_\delta$ above that theorem's rank threshold and so that
$\log R_\delta\ge\max\{100,1/a_\delta\}$.
For $\rho\ge R_\delta$,
\[
 \Phi(\rho)=\frac{\rho}{\log^2\rho}
 \le a_\delta\frac{\rho}{\log\rho}.
\]
The single-component condition also holds, since $1\le\log\rho/100$.
Theorem~\ref{thm:research-marked-size} applies.
\end{proof}

We also use the manuscript's small-cut estimate in the following form.
If $q$ disjoint cyclic regions avoid $V(W)$ and each has boundary at most an integer $D\ge1$, then
\begin{equation}\label{eq:general-small-cut-input}
 p(G,W)\le\frac12+\frac{(D+1)2^{D-1}}{q}.
\end{equation}
This statement permits loops and parallel edges.
It follows directly from the cut formula, the exceptional-partner bound, and the second-moment estimate.

\begin{lemma}[Localization]\label{lem:general-localization}
Let $G$ be a connected multigraph, and let $R$ be a nonempty connected proper vertex set.
Suppose that $G-R$ is connected, and write $d=|\delta_G R|\ge1$.
For each boundary edge $f$, with endpoint $z_f\in R$, add a distinct new vertex $t_f$ and an edge $z_ft_f$.
Add a cycle $W_R$ through the vertices $t_f$, padding it with new vertices if $d<3$.
Together with $G[R]$, this gives a graph $G_R$ satisfying
\[
 |E(W_R)|=\max\{d,3\},\qquad
 \beta(G_R)=\beta(G[R])+d.
\]
For uniform even subgraphs,
\begin{equation}\label{eq:general-localization-law}
 \Pr\bigl(\xi_G\cap E(G[R])\text{ is a forest}\bigr)=p(G_R,W_R).
\end{equation}
If $R\cap V(W)=\varnothing$, then $p(G,W)\le p(G_R,W_R)$.
\end{lemma}

\begin{proof}
Restriction outside an edge set $T$ maps the cycle space onto the cycle space after contracting every component of $T$.
For surjectivity, prescribe an even edge set in the contraction.
Its boundary inside each component of $T$ has even size.
A join in a spanning tree of that component completes the prescribed edge set.
This is a surjective linear map and therefore preserves uniformity on its image.

Contract $E(G-R)$ in $G$.
The resulting graph consists of $G[R]$ and one vertex receiving its $d$ boundary edges.
Contract $E(W_R)$ in $G_R$.
The same graph results, so the restrictions to $E(G[R])$ have the same law.
The added spokes have distinct leaf endpoints outside $W_R$.
They cannot create a cycle outside $W_R$, which proves \eqref{eq:general-localization-law}.
The stated inequality follows because $E(G[R])$ avoids $E(W)$.

For $d\ge3$, the construction adds $d$ vertices and $2d$ edges, giving the rank formula.
For $d=1,2$, padding the cycle adds equally many further edges and vertices.
Thus the rank formula remains valid.
\end{proof}

\begin{lemma}[Connected complementary shore]\label{lem:general-shore-closure}
Let $G$ and $W$ be connected, and let $R$ be a connected vertex set avoiding $V(W)$.
Let $C_0$ be the component of $G-R$ containing $W$, and put $R^+=V(G)\setminus V(C_0)$.
Then $R^+$ is connected, contains $R$, avoids $V(W)$, and satisfies
$\delta_G R^+\subseteq\delta_G R$.
Its complement is connected.
\end{lemma}

\begin{proof}
Every component of $G-R$ other than $C_0$ has an edge to $R$, since $G$ is connected.
Their union with $R$ is therefore connected.
Edges between $R^+$ and $C_0$ must have their other endpoint in $R$.
This gives the boundary inclusion and the remaining conclusions.
\end{proof}

\subsection{A shortest path in a cubic expansion}

\begin{lemma}\label{lem:general-shortest-reduction}
Suppose $P_{uv}(N)>2^{b-1}$.
There is a connected loopless multigraph $\Gamma$, with distinct terminals $u',v'$ and a shortest terminal path $A$, such that
\begin{enumerate}
\item $\beta(\Gamma)=b$ and $P_{u'v'}(\Gamma)\ge P_{uv}(N)$;
\item every nonterminal vertex of $\Gamma$ has degree three, and both terminals have degree one or two;
\item the length $\ell_0$ of $A$ is at most the length of every terminal path in $N$;
\item $|V(\Gamma)|\ge2b$, and $U=V(\Gamma)\setminus V(A)$ has size at least $2b-\ell_0-1$.
\end{enumerate}
Every vertex of $A$ has at most one neighbor incidence outside $A$.
\end{lemma}

\begin{proof}
By Proposition~\ref{prop:research-offroute}, no cyclic block lies outside the terminal route.
Delete the trees outside that route.
Suppress every degree-two vertex other than $u,v$.
This preserves rank and path count, and it can only shorten terminal paths.
No loop arises, since a cyclic block on the route retains its two distinct route terminals.
Call the resulting kernel $N_1$.
Every nonterminal vertex of $N_1$ has degree at least three.

Choose a shortest terminal path $A_1$ in $N_1$.
At each internal vertex of $A_1$, retain one vertex incident with its two path edges.
Attach all other incident edges through a rooted trivalent tree in its third branch.
At each terminal, retain one vertex incident with the path edge and attach its other edges through one rooted tree.
A single remaining incident edge needs no internal tree vertex.
Replace every other vertex of degree at least four by a trivalent tree.
No embedding restriction is imposed on these trees.

These replacements preserve rank.
Every original terminal path has a unique simple lift, so path count cannot decrease.
The image $A$ of $A_1$ has exactly the same length $\ell_0$.
Contracting the replacement trees maps any new terminal path to a terminal walk in $N_1$ of no greater length.
Every such walk has length at least the terminal distance in $N_1$, namely $\ell_0$.
Thus $A$ remains shortest.
This argument uses a shortest path after suppression, not an arbitrary path in the original graph.

All nonterminal degrees are three, and both terminal degrees are at most two.
If $n=|V(\Gamma)|$, the degree sum gives
\[
 2(n+b-1)=3(n-2)+\deg(u')+\deg(v').
\]
Hence $n=2b+4-\deg(u')-\deg(v')\ge2b$.
Subtracting $|V(A)|=\ell_0+1$ proves the last assertion.
The construction also gives the asserted bound on incidences from $A$ to $U$.
\end{proof}

\subsection{Wide distance layers give a region with small boundary}

Write $A=w_0w_1\cdots w_{\ell_0}$.
Let $L_j$ be the distance-$j$ layer from $u'=w_0$ in $\Gamma$.
Since $A$ is shortest, $w_j\in L_j$.
Put $U_j=L_j\cap U$ and $m_j=|U_j|$, taking $m_j=0$ for nonexistent layers.
For integers $i\le i'$, put
\[
 R(i,i')=U_i\cup\cdots\cup U_{i'},\qquad
 m(i,i')=\sum_{j=i}^{i'}m_j.
\]
Every edge joins vertices in the same or consecutive distance layers.
At most one edge joins each marked vertex to $U$.
Consequently,
\begin{equation}\label{eq:general-window-boundary}
 |\delta_\Gamma R(i,i')|
 \le3m_{i-1}+3m_{i'+1}+(i'-i+3).
\end{equation}
The last term counts the possible marked vertices in layers $i-1,\ldots,i'+1$.

For a parameter $\Theta\ge2^{36}$, call $R\subseteq U$ good if
\[
 |R|\ge\Theta,\qquad |\delta_\Gamma R|\le\Phi(|R|)/4.
\]
A good set need not be connected.

\begin{lemma}[Growing a window]\label{lem:general-window}
If $m_j\ge\Theta$ for some $j$, then a good set exists.
\end{lemma}

\begin{proof}
Start with the window $[j,j]$.
While a neighboring layer has mass greater than $\Phi(m(i,i'))/48$, include that layer.
Repeat until neither neighboring layer satisfies the condition.
Only finitely many layers are nonempty, so the procedure stops.

Let $\mu$ be the final mass, and set $L=\log\mu\ge36$.
Every step multiplies the current mass by more than $1+1/(48L^2)$.
Thus the number of steps is at most
\[
 \frac{\ln\mu}{\ln(1+1/(48L^2))}
 \le96L^2\ln\mu<67L^3.
\]
At termination, each neighboring layer has mass at most $\Phi(\mu)/48$.
Equation~\eqref{eq:general-window-boundary} therefore gives
\[
 |\delta_\Gamma R(i,i')|
 \le\Phi(\mu)/8+67L^3+3\le\Phi(\mu)/4.
\]
The final inequality follows from
$2^L\ge8L^2(67L^3+3)$ for $L\ge36$.
It holds at $L=36$, and the ratio of its left side to its right side increases thereafter.
Since $\mu\ge\Theta$, the final window is good.
\end{proof}

\begin{lemma}[No good set forces a long path]\label{lem:general-layer-count}
If no good subset of $U$ exists, then
\[
 \ell_0>\frac{2b}{\Theta+1}-(\Theta+1).
\]
\end{lemma}

\begin{proof}
Lemma~\ref{lem:general-window} gives $m_j<\Theta$ in every layer.
Let $R$ be the union of the layers $U_j$ with $j\ge\ell_0+2$.
No vertex of $A$ has a neighbor in $R$.
Therefore $|\delta_\Gamma R|\le3m_{\ell_0+1}<3\Theta$.
If $|R|\ge\Theta^2$, monotonicity of $\Phi$ and $\Phi(\Theta^2)/4\ge3\Theta$ would make $R$ good.
Hence $|R|<\Theta^2$.
Counting the remaining layers gives
\[
 2b-\ell_0-1\le|U|<(\ell_0+2)\Theta+\Theta^2.
\]
Rearranging proves the assertion.
\end{proof}

\subsection{A good set contradicts excess path density}

Add a new vertex $x$ and edges $u'x,xv'$ to $\Gamma$.
Call the resulting graph $H$, and mark the cycle $W=A\cup\{u'x,xv'\}$.
The proof of Lemma~\ref{lem:research-path-injection} also applies to multigraphs, giving
\begin{equation}\label{eq:general-path-forest}
 P_{u'v'}(\Gamma)\le p(H,W)2^b.
\end{equation}

\begin{lemma}[The good-set dichotomy]\label{lem:general-good-set}
Use $\delta,R_\delta,M,\Theta$ from Theorem~\ref{thm:general-linear-path-stability}.
If a good subset of $U$ exists, then $p(H,W)\le1/2+\delta$.
\end{lemma}

\begin{proof}
Let $R$ be good, with $m=|R|$ and $d=|\delta_\Gamma R|\le m/(4\log^2m)$.
Let $R_s$ be the vertex sets of the components of $\Gamma[R]$.
Their boundaries are disjoint and have total size $d$.
Every vertex in $R$ is cubic in both $\Gamma$ and $H$, so
\[
 \beta(\Gamma[R_s])=\frac{|R_s|-|\delta_\Gamma R_s|}{2}+1.
\]
Separate the components into four classes:
\begin{enumerate}
\item $|R_s|\ge M$ and $|\delta_\Gamma R_s|\le\Phi(|R_s|)/2$;
\item $|R_s|\ge M$ and $|\delta_\Gamma R_s|>\Phi(|R_s|)/2$;
\item $|R_s|<M$ and $\Gamma[R_s]$ is a tree;
\item $|R_s|<M$ and $\Gamma[R_s]$ contains a cycle.
\end{enumerate}

Class 2 has total size less than $2d\log^2m\le m/2$.
Indeed, each of its components satisfies
$|R_s|<2|\delta_\Gamma R_s|\log^2|R_s|\le2|\delta_\Gamma R_s|\log^2m$.
A tree component satisfies $|\delta_\Gamma R_s|=|R_s|+2$.
Thus class 3 has total size less than $d\le m/4$.
A cyclic component in class 4 has boundary at most its size, hence less than $M$.

Suppose first that class 1 is empty.
Class 4 then has total size greater than $m/4$, so it contains more than $m/(4M)$ components.
These are disjoint cyclic regions avoiding $W$, with boundaries less than $M$.
By \eqref{eq:general-small-cut-input},
\[
 p(H,W)\le\frac12+
 \frac{(M+1)2^{M-1}\,4M}{m}
 \le\frac12+\delta.
\]
The last inequality uses $m\ge\Theta$.

Suppose instead that class 1 contains $R_s$.
Apply Lemma~\ref{lem:general-shore-closure} in $H$ to obtain $R^+\supseteq R_s$ with connected complement.
It avoids $V(W)$, so it is contained in $U$ and all its vertices are cubic.
Write $d^+=|\delta_H R^+|$ and
\[
 \rho=\beta(H[R^+])+d^+
 =\frac{|R^+|+d^+}{2}+1
 \ge|R_s|/2\ge R_\delta.
\]
Since $\delta_HR_s=\delta_\Gamma R_s$, the boundary comparison gives
\[
 d^+\le|\delta_\Gamma R_s|
 \le\Phi(|R_s|)/2
 \le\Phi(|R_s|/2)
 \le\Phi(\rho).
\]
The middle inequality uses $\log(s/2)=\log s-1$.

Localize at $R^+$ using Lemma~\ref{lem:general-localization}.
If $d^+\ge3$, its marked cycle has at most $\Phi(\rho)$ edges by the last display.
If $d^+=1$ or $2$, its marked cycle has three edges, which still meets that bound since $\rho\ge2^{100}$.
Subdivide any parallel interior edges until the localized graph is simple.
Such subdivisions preserve rank, marked-cycle length, and outside-forest probability.
Lemma~\ref{lem:general-fe2} now yields
\[
 p(H,W)\le p(H_{R^+},W_{R^+})\le1/2+\delta.
\]
This completes both cases.
\end{proof}

\begin{proof}[Proof of Theorem~\ref{thm:general-linear-path-stability}]
Use Lemma~\ref{lem:general-shortest-reduction}.
The density assumption and $\delta=\varepsilon/2$ give
\[
 p(H,W)\ge\frac{P_{u'v'}(\Gamma)}{2^b}
 \ge\frac12+\varepsilon>\frac12+\delta.
\]
Lemma~\ref{lem:general-good-set} excludes every good subset of $U$.
Lemma~\ref{lem:general-layer-count} proves \eqref{eq:general-linear-strong} for $\ell_0$.
Every original terminal path has length at least $\ell_0$, so the same inequality holds for its length $\ell$.

Put $Q=\Theta+1$.
If $b\ge Q^2$, then $\ell>2b/Q-Q\ge b/Q\ge b/Q^2$.
If $b<Q^2$, then $\ell\ge1>b/Q^2$.
Thus $c_\varepsilon=Q^{-2}$ works at every rank.
\end{proof}

\subsection{Structural consequences}

\begin{corollary}\label{cor:general-short-path-density}
For any sequence of connected graphs with ranks $b\to\infty$ and terminal distances $o(b)$,
\[
 P_{uv}(N)\le(1/2+o(1))2^b.
\]
If $P_{uv}(N)\ge(1-\alpha)2^b$ for a fixed $0\le\alpha<1/2$, every terminal path has length
$\Omega_\alpha(b)$.
\end{corollary}

\begin{proof}
A fixed positive density gap in the first assertion would contradict Theorem~\ref{thm:general-linear-path-stability}.
For the second assertion, use $\varepsilon=1/2-\alpha$.
\end{proof}

\begin{corollary}[Bounded width of the terminal kernel]\label{cor:general-bounded-pathwidth}
Under the hypotheses of Theorem~\ref{thm:general-linear-path-stability}, remove all blocks outside the terminal route and suppress nonterminal degree-two vertices.
The resulting terminal kernel has pathwidth at most $2\lceil\Theta\rceil+1$.
The same bound holds for the cubic expansion in Lemma~\ref{lem:general-shortest-reduction}.
\end{corollary}

\begin{proof}
The proof excludes good sets, so Lemma~\ref{lem:general-window} gives $|U_j|<\Theta$ in every distance layer.
At most one marked path vertex lies in each layer.
Therefore each full layer has at most $\lceil\Theta\rceil+1$ vertices.
The bags $L_j\cup L_{j+1}$ form a path decomposition, since every edge joins equal or consecutive layers.
Their sizes give the stated pathwidth bound.
The terminal kernel is obtained from the cubic expansion by contracting its replacement trees.
Pathwidth does not increase under taking minors.
\end{proof}

The restriction to the terminal kernel matters.
Trees attached outside the terminal route can have arbitrarily large pathwidth without changing rank or terminal path count.

\section{General additive path stability}
\label{sec:general-additive-path-stability}

The bounded-layer conclusion of the preceding proof yields a stronger result.
It forces every terminal path to have length at least the rank, up to a constant depending on the density gap.
This deduction was developed by the Codex path-stability agent and independently audited by the root and four-regular agents.
It uses Fable's localization and layer argument through the preceding section.
The result has not been formalized in Lean.

\begin{theorem}[Additive path stability]\label{thm:general-additive-path-stability}
For every $0<\varepsilon\le1/2$, there is a constant $C_\varepsilon$ such that the following holds.
If $N$ is a connected simple graph of rank $b$, with distinct terminals $u,v$, and
\[
 P_{uv}(N)\ge(1/2+\varepsilon)2^b,
\]
then every $u$--$v$ path has length at least
\[
 b-C_\varepsilon.
\]
In particular, $\operatorname{dist}_N(u,v)\ge b-C_\varepsilon$.

Let $\Theta$ be the constant from Theorem~\ref{thm:general-linear-path-stability}, and put
\[
 B=\lceil\Theta\rceil,\qquad L=20B+2.
\]
One may take
\begin{equation}\label{eq:general-additive-constant}
 C_\varepsilon=\frac{2L(L+1)2^{L-1}}{\varepsilon}.
\end{equation}
For the cubic expansion $\Gamma$ and shortest path $A$ from Lemma~\ref{lem:general-shortest-reduction},
one also has $\beta(\Gamma-V(A))\le C_\varepsilon$.
\end{theorem}

The size of this constant is not optimized.
Its finiteness gives the qualitative additive conclusion.

\subsection{A short cycle in any cyclic layer interval}

Retain $\Gamma,A,U_j$ from the preceding proof.
All vertices outside $A$ are cubic, and each vertex of $A$ has at most one incidence outside $A$.
Assume that $|U_j|\le B$ in every distance layer.

\begin{lemma}[Minimum layer span]\label{lem:additive-minimum-span}
If the subgraph induced by $\bigcup_{j\in I}U_j$ contains a cycle for an interval $I$ of layer indices,
then it contains a cycle of length at most $20B+2$.
\end{lemma}

\begin{proof}
Among all cycles in that subgraph, choose one whose layer span is smallest.
Let its minimum and maximum layer indices be $a,b'$, and put $s=b'-a$.
Let $R=U_a\cup\cdots\cup U_{b'}$.
Every cycle in $\Gamma[R]$ has span at least $s$, by the choice of the first cycle.
Since it lies in this same interval, it must meet $U_a$.
Thus $\Gamma[R]-U_a$ is a forest.

Deleting a vertex of degree at most three decreases rank by at most two.
Consequently,
\[
 \beta(\Gamma[R])\le2|U_a|\le2B.
\]
Put $m=|R|$, $d=|\delta_\Gamma R|$, and $c=\operatorname{comp}(\Gamma[R])\ge1$.
Cubicity outside $A$ gives the exact identity
\[
 m=2\beta(\Gamma[R])+d-2c.
\]
The layer-boundary estimate \eqref{eq:general-window-boundary} yields
\[
 d\le6B+s+3.
\]
Combining these bounds gives
\begin{equation}\label{eq:additive-span-size}
 m\le10B+s+1.
\end{equation}

The chosen cycle has length at least $2s$.
Indeed, its two arcs between a minimum-layer vertex and a maximum-layer vertex each have at least $s$ edges.
This also holds when $s=0$.
Hence $2s\le m\le10B+s+1$, so $s\le10B+1$.
Equation~\eqref{eq:additive-span-size} now gives $m\le20B+2$.
The chosen cycle has at most $m$ edges, proving the assertion.
\end{proof}

\begin{lemma}[Greedy deletion of whole windows]\label{lem:additive-layer-deletion}
Put $J=\Gamma-V(A)$ and $L=20B+2$.
There are $k$ vertex-disjoint cycles in $J$, each of length at most $L$, such that
\[
 \beta(J)\le2Lk.
\]
\end{lemma}

\begin{proof}
Start with all unmarked layers.
If one remaining layer interval contains a cycle, choose a cycle of minimum span in that interval.
Delete its entire tight layer window $R$, as defined in Lemma~\ref{lem:additive-minimum-span}.
That lemma's proof gives $|R|\le L$.
The remaining vertices again form whole layer intervals, and no edge joins two different intervals.
Repeat until the remaining graph is a forest.

Suppose that this procedure selects $k$ cycles, and let $Z$ be the union of the deleted windows.
The selected cycles are vertex-disjoint, have length at most $L$, and satisfy $|Z|\le Lk$.
Since $J$ has maximum degree at most three, deleting $Z$ decreases rank by at most $2|Z|$.
Therefore $\beta(J)\le2|Z|\le2Lk$.
\end{proof}

\subsection{Completing the additive estimate}

\begin{proof}[Proof of Theorem~\ref{thm:general-additive-path-stability}]
Use Lemma~\ref{lem:general-shortest-reduction} to obtain $\Gamma$ and its shortest terminal path $A$ of length $\ell_0$.
In the proof of Theorem~\ref{thm:general-linear-path-stability}, the density assumption excludes every good set.
Lemma~\ref{lem:general-window} therefore gives $|U_j|<\Theta\le B$ for every layer.

Use Lemma~\ref{lem:additive-layer-deletion} to choose $k$ vertex-disjoint cycles of length at most $L$ in $J=\Gamma-V(A)$.
Each of their vertex sets is a cyclic region with boundary at most $L$ in $\Gamma$.
Every vertex of such a cycle has degree three and uses two incidences in that cycle.
Add the new two-edge terminal path and mark its union with $A$, giving $H,W$ as before.
The chosen region boundaries are unchanged, and all the regions avoid $V(W)$.

If $k\ge1$, the path injection and \eqref{eq:general-small-cut-input} give
\[
 \frac12+\varepsilon\le p(H,W)
 \le\frac12+\frac{(L+1)2^{L-1}}{k}.
\]
Thus $k\le(L+1)2^{L-1}/\varepsilon$.
The same upper bound holds when $k=0$.
Lemma~\ref{lem:additive-layer-deletion} yields
\[
 \beta(J)\le2Lk\le C_\varepsilon.
\]

It remains to compare this rank with $b-\ell_0$.
The degree sum at the $\ell_0+1$ vertices of $A$ is at most $3(\ell_0-1)+4=3\ell_0+1$.
These vertices contain the $\ell_0$ path edges.
Deleting them therefore removes at most $2\ell_0+1$ edges.
If $J$ is nonempty, its rank satisfies
\[
 \beta(J)\ge b+(\ell_0+1)-(2\ell_0+1)=b-\ell_0.
\]
If $J$ is empty, then $|E(\Gamma)|\le(3\ell_0+1)/2$, so
\[
 b=|E(\Gamma)|-\ell_0\le(\ell_0+1)/2\le\ell_0.
\]
Thus $b-\ell_0\le\beta(J)$ in both cases.
It follows that $\ell_0\ge b-C_\varepsilon$.
Every original terminal path has length at least $\ell_0$, which proves the theorem.
\end{proof}

\begin{corollary}\label{cor:general-rank-deficit}
For any sequence of connected graphs and terminal pairs with
$\beta(N)-\operatorname{dist}_N(u,v)\to\infty$,
\[
 P_{uv}(N)\le(1/2+o(1))2^{\beta(N)}.
\]
\end{corollary}

\begin{proof}
A fixed positive density gap would bound the rank-distance deficit by Theorem~\ref{thm:general-additive-path-stability}.
\end{proof}

The conclusion does not yet give an exponential error in the rank-distance deficit.
Such a rate would require quantitative control beyond the constants supplied here.

The script \texttt{check\_additive\_layer\_structure.py} checks the finite layer argument on 157 explicit cubic expansions.
It verifies 301 selected windows, including examples with 322 vertices and outside rank 81.
The checks cover the minimum-span feedback set, cubic rank identity, boundary bound, and final greedy rank estimate.
They support these finite identities and do not replace the probability argument.

\end{document}
